\documentclass[10pt,a4paper]{amsart}
\usepackage[margin=19mm]{geometry}
\usepackage{amsmath,amssymb,mathtools}
\usepackage[scr=rsfs,cal=euler]{mathalfa}
\usepackage{xfrac}
\usepackage[unicode,hidelinks,bookmarks=false]{hyperref}
\newcommand{\ie}{i.e.\ }
\newcommand{\cf}{cf.\ }
\newcommand{\resp}{respectively\ }
\newcommand{\Subsection}{\subsection}
\providecommand{\nobreakdash}{\nobreak\hbox{-}}
\setlength{\emergencystretch}{2em}
\multlinegap0pt
\usepackage{bm}
\usepackage{bbm}
\usepackage{dsfont}
\usepackage{xspace}
\usepackage{cancel}
\usepackage{mathtools}
\usepackage{amsthm}
\makeatletter
\@ifundefined{prop}{\newtheorem{prop}{Proposition}}{}
\makeatother
\title[PDE-Free Mass-Constrained Learning of Complex Systems with H]{PDE-Free Mass-Constrained Learning of Complex Systems with Hidden States}
\author{Gianmaria Viola, Alessandro Della Pia, Lucia Russo, Ioannis Kevrekidis, Constantinos Siettos}
\date{Source: arXiv:2510.17657v3 (CC BY 4.0)}
\begin{document}
\maketitle

\noindent\emph{Source and licence.} PDE-Free Mass-Constrained Learning of Complex Systems with Hidden States. Version arXiv:2510.17657v3 (CC BY 4.0), distributed under the Creative Commons Attribution 4.0 International licence, as recorded in the arXiv licence metadata for this exact version and shown on the abstract page https://arxiv.org/abs/2510.17657v3. Full rights notice: licenses/arxiv-2510.17657-CC-BY-4.0.txt.\par\smallskip

\noindent\textit{Editorial scope.} Counted unit: the Prop whose source label is \texttt{prop1} in the article (source0.tex lines 192--218, source label \texttt{prop1}), delivered together with the article's own complete original proof of it (source0.tex lines 220--245, one \texttt{proof} environment of 26 lines). No mathematics is written, repaired or re-derived: statement and proof are the article's own text line for line. Required context retained uncounted: none. Every definition, display and label of those lines is retained unchanged. Omitted from this package and not redistributed: 1--191, 219--219, 246--1066. Nothing inside a retained excerpt is omitted or altered: every retained body is delivered exactly as the article's lines are written, so no omission receipt of any kind accompanies this package. Citations: the retained excerpts carry no citation command at all, so no bibliography entry of the article is transcribed and no reference is redistributed. Reference adaptation: none was needed and none was made; every reference inside the delivered unit resolves inside it, so no unresolved reference or citation is produced. Printed slips kept literal: no printed word, symbol, hypothesis or number of the counted statement or of its proof was altered. Layout and class: the article's own class is \texttt{article} with packages including fix-cm, longtable, arxiv, lipsum, placeins, algpseudocodex, float, babel, inputenc, fontenc, hyperref, url, booktabs, stackrel; the wrapper is the lane's pinned \texttt{amsart} editorial wrapper at 10pt on A4, which restates the theorem-like environments the retained text needs and adds the article's own macro definitions that the retained text needs (none), with extra wrapper packages (bm, bbm, dsfont, xspace, cancel, mathtools). The delivered numbering is the unit's own. Assets: no retained span contains a figure environment, a graphics-inclusion command, a PDF-page inclusion, a table or a TikZ picture; no image file is copied into this package, the package declares no asset dependency and no image file is redistributed with it. Licence: version arXiv:2510.17657v3 of this article is distributed under the Creative Commons Attribution 4.0 International licence, as recorded in the arXiv licence metadata for this exact version and shown on the abstract page https://arxiv.org/abs/2510.17657v3. A full rights notice is delivered at licenses/arxiv-2510.17657-CC-BY-4.0.txt. Characters: every retained excerpt is pure ASCII, so the delivered engine, XeLaTeX with its default Latin Modern text font, reports no missing character.
\begin{prop}
\label{prop1}
Let the dataset be 
\begin{equation}
    \mathbf{X} = [\mathbf{x}_1, \mathbf{x}_2, \dots, \mathbf{x}_M] \in \mathbb{R}^{N \times M},
\end{equation}
with columns $\mathbf{x}_k \in \mathbb{R}^{N}$ representing normalized to sum to 1 density fields. If the total mass of the original data is preserved, i.e.,
\begin{align}
    \begin{bmatrix}
      \sum_{i=1}^{N} x_{i1} & \sum_{i=1}^{N} x_{i2} & \dots & \sum_{i=1}^{N} x_{iM}
    \end{bmatrix}
    = \mathbf{1}_M^\top,
\end{align}
then the mass of the \textit{reconstructed/decoded} field computed by POD as
\begin{equation}
\hat{\mathbf{X}} = \mathbf{U}_d 
\mathbf{U}_d^\top \tilde{\mathbf{X}} +  \mathbf{X} (\mathbf{I}_M - \mathbf{H}), \quad \mathbf{H} = \mathbf{I}_M - \frac{1}{M} \mathbf{1}_M \mathbf{1}_M^\top,
\label{eq:recon_theorem}
\end{equation}
where $\mathbf{U}_d = [\mathbf{u}_1,\mathbf{u}_2,\dots, \mathbf{u}_d] \in \mathbb{R}^{N \times d}$ is the orthonormal basis formed by the first $d$ left singular vectors of the SVD of the centered matrix $\tilde{\mathbf{X}} = \mathbf{X} \mathbf{H}$, is also preserved, i.e.,
\begin{align}
    \begin{bmatrix}
      \sum_{i=1}^{N} \hat{x}_{i1} & \sum_{i=1}^{N} \hat{x}_{i2} & \dots & \sum_{i=1}^{N} \hat{x}_{iM}
    \end{bmatrix}
    = \mathbf{1}_M^\top.
\end{align}
\end{prop}

\begin{proof}
Since the sum of each column of $\mathbf{X}$ is 1 by hypothesis, we have
\[
\mathbf{1}_N^\top \mathbf{X} = \mathbf{1}_M^\top.
\]
Then, for the centered matrix $\tilde{\mathbf{X}} = \mathbf{X} \mathbf{H}$, we obtain
\begin{align}
\mathbf{1}_N^\top \tilde{\mathbf{X}} = \mathbf{1}_N^\top \mathbf{X} \mathbf{H} = \mathbf{1}_M^\top \mathbf{H} = \mathbf{1}_M^\top - \frac{1}{M} \mathbf{1}_M^\top \mathbf{1}_M \mathbf{1}_M^\top = \mathbf{1}_M^\top - \mathbf{1}_M^\top = \mathbf{0}_M^\top.
\end{align}
This implies that $\mathbf{1}_{N}$ is orthogonal to the column space of $\tilde{\mathbf{X}}$, and therefore to the left singular vectors of its SVD, which provide an orthogonal basis for that column space. Thus,
\begin{equation}
    \mathbf{1}_{N}^\top \mathbf{U}_d = \boldsymbol{0}^\top_d,
\end{equation}
where $d=1,2,\dots, r$, and $r$ is the rank of $\mathbf{X}$. Multiplying Eq.~\eqref{eq:recon_theorem} by $\mathbf{1}_{N}^\top$ yields
\begin{equation}
\mathbf{1}_N^\top \hat{\mathbf{X}} = \mathbf{1}_N^\top \mathbf{U}_d \mathbf{U}_d^\top \tilde{\mathbf{X}} + \mathbf{1}_N^\top \mathbf{X} (\mathbf{I}_M - \mathbf{H})  = \mathbf{0}_M^\top + \frac{1}{M} \mathbf{1}_N^\top \mathbf{X} \mathbf{1}_M \mathbf{1}_M^\top = \mathbf{1}_M^\top,
\end{equation}
since $\mathbf{1}_M^\top \mathbf{1}_M = M$. We therefore obtain
\begin{align}
    \mathbf{1}_{N}^\top \hat{\mathbf{X}} =
    \begin{bmatrix}
    \sum_{i=1}^{N} \hat{x}_{i1} &  \sum_{i=1}^{N} \hat{x}_{i2} & \dots & \sum_{i=1}^{N} \hat{x}_{iM}  
    \end{bmatrix} = \mathbf{1}^\top_{M},
\end{align}
hence showing that the total mass in each snapshot is preserved under POD reconstruction.
\end{proof}

\end{document}
